\documentclass[letterpaper,12pt,twoside]{book}  % book class
%%s\usepackage[spanish,es-nodecimaldot,es-tabla]{babel} % language/hyphenation
\usepackage[utf8]{inputenc}  % input encoding
\usepackage{graphicx} % graphics library (REQUIRED by CIMATpreamble)
\usepackage[svgnames,table]{xcolor}  % color library (REQUIRED by CIMATpreamble)
\usepackage{CIMATpreamble}  % definition of the CIMAT cover page
\usepackage{mypreamble} % custom preamble
\graphicspath{{img/}}  % set path for images folder

%-----------------------------------------------------------------------------%
% Coverpage template requested by CIMAT for the master thesis in LaTeX version 

%-----------------------------------------------------------------------------%
% README: Instructions & Comments
%
% > The package 'CIMATpreamble' is required for the CIMAT cover page and 
%   it depends on the declared libraries above
% > To use the COVER template only fill the Cover/Titlepage block 
%   with your information and make sure the '\maketitle' command is enabled
%
% > The rest of the LaTeX code was built according to the guideline provided by 
%   the ShareLatex/Overleaf team:
%   "How to Write a Thesis in LaTeX (Part 1): Basic Structure"; you will find  
%   the custom settings (link colors, caption specs., etc.) in the 
%   'mypreamble.sty' file, so you are free to use or modify that code 
% > Much of the code for the cover page is based on the Overleaf template    
%   created by Salvador Pedraza Espitia: 
%   https://www.overleaf.com/latex/templates/unam-dissert/mypjkyrmhgns
% > The essential idea of this code is to provide a minimal and functional  
%   LaTeX document, allowing us to focus on the content of our thesis work;
%   I hope it be useful for you :) !
% 
% > Credits: Kenny Yahir Mendez Ramirez, Cristal E. Mendez Ramirez
% > For comments and suggestions please send a e-mail to kenny.yahir.mz@gmail.com

%-----------------------------------------------------------------------------%
% Cover/Title page: filzzxl out with your personal information

\author{Gustavo Hernández Angeles}
\documentType{T E S I S}  % change to "T E S I N A" if it is the case
\title{La Voz del Turista: Flujo de Trabajo Basado en Modelos de Lenguaje para la Mejora de la Ventaja Competitiva en Pueblos Mágicos}
\degree{Maestro en Cómputo Estadístico}
\supervisor{Dr. Norberto Alejandro Hernández Leandro}
%\supervisorSecond{Dra. M. F.}  % optional command
\cityandyear{Monterrey, Nuevo León, Xxxx de 2026}

%-----------------------------------------------------------------------------%
% Document body

\begin{document}

% Coverpage
\maketitle  % this enable the CIMAT title page!

\thispagestyle{empty}  % blank page after title page

\frontmatter

% Dedication
\chapter*{}
\begin{flushright}%
  \emph{To my family.} \\[2cm]
  

  \thispagestyle{empty}
\end{flushright}

% Acknowledgements
\chapter*{Acknowledgements}
\addcontentsline{toc}{chapter}{Acknowledgements}

\noindent To my parents Octavio and Abigail, who have provided me with immense and unconditional support. They are my inspiration and strength. \\[0.2cm]

\noindent To my brothers, Arael and Javier, for their constant support and for being role models. \\[0.2cm]

\noindent To my friends, for making this journey more bearable with laughter and constant support. \\[0.2cm]

\noindent To Yazmin, for the great support, patience, and continuous joys.

% Abstract
\chapter*{Abstract}
\addcontentsline{toc}{chapter}{Abstract}

The use of LLMs has revolutionized the way organizations can analyze and leverage large volumes of unstructured data. In particular, tourist destinations face the challenge of extracting valuable information from tourists' opinions and comments to improve their competitive advantage. Traditionally, opinion analysis requires significant manual effort or NLP techniques that may not fully capture the context and nuances of human language.

In this thesis, we present a workflow based on LLMs to analyze tourists' opinions in Mexico's \textit{Pueblos Mágicos}, aiming to identify insights in their hotel and restaurant markets, as well as natural and historical tourist resources, promoting continuous improvement of competitive advantage in Pueblos Mágicos. We propose a Map-Reduce workflow that integrates RAG for review retrieval using Sentence Transformer models (also known as bi-encoders) that capture the semantics of opinions. Subsequently, LLMs are employed for analyzing and synthesizing the retrieved information, generating executive reports that highlight key areas for improvement and strategic opportunities.


% Keywords
\paragraph{Keywords:} Large Language Models, Competitive Advantage, Pueblos Mágicos, Tourism, Opinion Analysis, Retrieval-Augmented Generation, Sentence Transformers.



% Table of contents and list of figures
\tableofcontents

\listoffigures

% Chapters
\mainmatter

\spacing{1.5}  % interline  (double) space

% put the content of your chapters in the .tex files located in the 'chapters/' folder; 
% the use of '\include' instead of '\input' command could enhance the  compilation process, see the reference below

\chapter{Introducción}

En las últimas décadas se ha presenciado un aumento exponencial en la creación e intercambio de información por los usuarios en plataformas digitales, fenómeno conocido como \textit{electronic Word of Mouth} (eWOM). Esta dinámica ha transformado la manera en que las organizaciones comprenden y responden a las necesidades y preferencias de sus clientes respecto a un producto o servicio concreto \citep{mishra2016ewom}. En la industria hotelera, se han realizado numerosos estudios que evidencian el impacto significativo del eWOM en la percepción de los consumidores y su influencia en la toma de decisiones (tanto de los potenciales usuarios como de los propios proveedores de servicios) \citep{cantallops2014new}. Un comportamiento similar se observa en el sector restaurantero, donde las opiniones y reseñas son cruciales para conocer la percepción de los clientes sobre la calidad de los alimentos, el servicio y la experiencia general con el fin de desarrollar estrategias de administración efectivas \citep{bilgihan2018identifying, ong2012perceived}.



\section{Hipótesis de trabajo} \label{sec:hypothesis}

\section{Objetivos} \label{sec:objective}

\section{Contribución} \label{sec:contribution}

En literatura existe una basta colección de trabajos que abordan el problema...

\section{Organización de la tesis}








\chapter{Marco teórico} \label{chap:background}


\chapter{Análisis exploratorio}\label{chap:eda}


\chapter{Diseño experimental y análisis de resultados} \label{chap:design_results}

\section{Diseño experimental} \label{sec:experimental_design}

\section{Análisis de resultados} \label{sec:results_analysis}

\include{chapters/conclusions}

% Bibliography
\bibliographystyle{apacite}  % APA style according to apacite package: see mypreamble.sty line 21
\bibliography{references}

% Appendix
\appendix
\include{chapters/appendix}

\backmatter

\end{document}

% LaTeX references (technical issues, workarounds)

% https://www.overleaf.com/learn/latex/Management_in_a_large_project ('including files' section)
% https://tex.stackexchange.com/questions/64839/how-to-change-listing-caption
% https://stackoverflow.com/questions/2709898/change-list-of-listings-text
% https://tex.stackexchange.com/questions/664/why-should-i-use-usepackaget1fontenc